% GENERATED by tools/preliminary_extraction_table.py --- do not hand-edit.
% Numbers: data/gap/archive/v0_4_extraction_run.json (computed from the v0.4 GAP
%   on archive/attacker-profiling and from data/gap/gap_v0.5.json).
% Wording: data/gap/archive/preliminary_extraction_labels.json.
% DRAFT STATE (2026-10-04 appendix pass): caption session-drafted, ratify on read.
% Requires booktabs (already in the preamble).

\begin{table}[htbp]
\centering
\caption[Edges between tactics found by mining and by attack flows]{Three ways of finding the edges between tactics (Section~\ref{sec:technique-graph}), run in one earlier build of the attack graph (2026-04-19) over the same ATT\&CK techniques. \emph{Edges between tactics} counts the ordered pairs of different tactics each found, of the 182 that ATT\&CK's 14 tactics then allowed.}
\label{tab:preliminary-extraction}
\tablestyle
\begin{tabular}{@{}P{0.13\textwidth} P{0.20\textwidth} P{0.14\textwidth} r P{0.26\textwidth}@{}}
\toprule
Way & An edge rests on & Direction from & \shortstack[r]{Edges between\\tactics} & Why it was not used \\
\midrule
Co-occurrence mining & the same APT group or tool uses both techniques & an assumed order of the tactics & 19 (10\,\%) & Using two techniques together says nothing about which came first, and an assumed order allows no edge back. \\
Keyword mining & one technique's ATT\&CK description names the other & an assumed order of the tactics & 25 (14\,\%) & A description that names another technique refers to it; it does not say the attacker needs it first. \\
\textbf{Attack flows} & an analyst drew one step after the other & the analyst's drawing & 91 (50\,\%) & Adopted. \\
\bottomrule
\end{tabular}
\end{table}
